\chapter{Edit in Rhino, and Bake}\label{ch:edit}

\section{Edit in Rhino}\label{sec:editinrhino}

Some fixes are easier to find by hand than with numbers: drag a control point until the curve looks
right, then read off what it should have been. \ui{Edit in Rhino} makes a paused script's value into
real Rhino objects for that. Change them with any Rhino tool, and the change is written back into the
paused script (\see{sec:writeback}). Python~3 scripts only.

\paragraph{Starting it.} Choose \ui{Edit in Rhino} in \vn{Watch}'s toolbar, in a tooltip, or in the
right-click menu of a name or a Variables row; or type \code{edit('crv')} in \vn{Python}. The value
becomes objects on the layer \file{Script Editor Watch::Edit}, one object per list item. Its drawing is
hidden meanwhile, and its name dot moves to the edited objects and with them.

\paragraph{Editing.} Move them, drag control points, use the Gumball, \code{Scale}, \code{Rebuild}, even
Undo. When a command ends, its edits are written back together, with one redraw. The new geometry is
turned back into the variable's type and assigned in the paused frame:

\begin{xltabular}{\linewidth}{@{}P{0.38\linewidth}L@{}}
\toprule
\textsf{\textbf{The variable was}} & \textsf{\textbf{Written back from}} \\ \midrule
\code{Point3d}                 & a point \\
\code{Line}, \code{Polyline}, \code{Circle}, \code{Arc} & a curve, if it still is one \\
\code{Box}                     & a Brep \\
\code{Vector3d}                & its arrow \\
\code{Curve}, \code{Brep}, \code{Mesh} & as they are \\
a list item                    & written as \code{curves[2] = …} \\
raw points, polylines (Interpret as) & written into the list; lists stay the same lists, tuples become tuples \\
your class with a visualizer   & through its \code{@rhinoviv.inverse} function (\cref{ch:visualizers}) \\
\bottomrule
\end{xltabular}

\paragraph{Undo.} \keys{Ctrl+Z} and Redo write the earlier or later geometry back too.

\paragraph{At every pause} the objects take the script's current values. \textbf{When the run ends}, they are
removed; \code{miniVivisection StopEditing} removes them earlier.

\paragraph{Limits.}
\begin{itemize}
  \item Write-back works only while paused. With the script on a thread of its own, it needs FreePause.
  \item An edit that changes the kind of geometry, a circle made into a free-form curve, say, is reported
        and not written back.
  \item Planes can't be edited.
  \item Other interpretations of raw data (Interpret as) are drawn only.
\end{itemize}

\section{Auto edit}\label{sec:autoedit}

With \ui{Auto edit} on (the default), every shown watch whose value is geometry becomes editable objects
at each pause, and right away for a watch added or shown while paused. It is \vn{Watch}'s toolbar switch,
and one setting for all editors (\setting{Auto edit}): every \vn{Watch} tab and the settings show the same
state. Turning it off removes those objects.

\section{Bake}\label{sec:bake}

\ui{Bake} copies drawn values into the document, on the layer \file{Script Editor Watch::Baked}. Each copy
is named after its expression and uses its color; a list's items are grouped. It is one undo step.

Bake from \vn{Watch}'s toolbar (the selected watch, or per thread row, \see{sec:everythread}), from a tooltip,
from the right-click menu (\ui{Bake to Rhino}), with \code{bake(value)} in \vn{Python}, or with
\code{miniVivisection BakeAll}, which bakes everything drawn right now.

\begin{tip}[Pause before a crash]
A script that crashes after a long computation loses its results. Set a breakpoint just before the line
that fails, and bake what was computed so far.
\end{tip}

\begin{lesson}{lessons 03, 04}
\lessonfile{03_bake_before_a_crash.py}: Bake. \lessonfile{04_edit_in_rhino.py}: Edit in Rhino, Auto edit.
\end{lesson}
